\section{Програмна реализация}
\label{sec:implementation}

Настоящата глава описва как проектните решения от Глава~\ref{sec:design} се материализират в конкретен програмен код. Изложението следва три равнища. Първото обхваща реализацията на асинхронния слой и неговите ключови компоненти --- \code{Task<T>}, обединението от работни нишки (\code{ThreadPool}), \code{run\_on\_pool()} и реакторния \code{IoContext}. Второто описва конкретните реализации на конекторите за основните системи за съхранение (backend): SQLite, PostgreSQL, MySQL, MongoDB и Redis. Третото равнище е верификацията: единични (unit), интеграционни (integration) и реално свързани с живи системи за съхранение тестове, заедно с емпиричната оценка на асинхронния модел.

\subsection{Методика и среда за реализация}
\label{subsec:method_environment}

Реализацията е изградена изцяло върху стандарта C++23~\cite{cpp23standard}, с използване на функционалност за корутини, въведена в C++20~\cite{cpp20standard,p0912r5,p1063r0}. Кодовата база е организирана по подсистеми. Заглавните файлове са разположени в директории, които съответстват на архитектурните слоеве: \path{lib/include/ORM/entity/} за обектния слой, \path{lib/include/ORM/query/} за слоя на заявките, \path{lib/include/ORM/connector/} за общата фасада и асинхронната обвивка, \path{lib/include/ORM/async/} за \code{Task<T>}, \code{ThreadPool} и \code{IoContext}, и \path{lib/include/ORM/db/connectors/} с отделни поддиректории за всяка конкретна система за съхранение. Системата за изграждане е CMake; конструирането и изпълнението на тестовете се извършва чрез \code{ctest}.

Поддръжката за живи (live) системи за съхранение е условна. Тестовете срещу реален PostgreSQL, MySQL, MongoDB, Redis, Neo4j и Cassandra инстанциран сървър се изпълняват само когато съответната среда е активирана чрез променливи на обкръжението и Docker-базирана инфраструктура. Това позволява стандартното изграждане да остане преносимо, докато по-задълбоченото верифициране се извършва селективно. За единичните тестове и базовия интеграционен сценарий този разделителен подход не е необходим --- те се компилират и изпълняват винаги.

\subsection{Реализация на асинхронния слой}
\label{subsec:async_implementation}

Асинхронният слой е разположен над синхронната фасада \code{db<DB>} и приема статично конструираното междинно представяне на заявката (query intermediate representation, query IR) като готов вход. Реализацията се фокусира върху три аспекта: жизнения цикъл на корутината (\code{Task<T>}), универсалния път на изнасяне (\code{run\_on\_pool()}) и реакторното обработване на готовността при драйверите с неблокиращи интерфейси (\code{IoContext}).

\subsubsection{\code{Task<T>}, продължение и завършване}

\code{Task<T>} е централната абстракция в \path{lib/include/ORM/async/task.hpp}. Имплементацията дефинира \code{promise\_type}, който поддържа: \code{initial\_suspend()}, връщащ \code{std::suspend\_always} (мързеливо стартиране); \code{final\_suspend()}, чрез който се възобновява запазеното продължение (continuation); \code{return\_value(T)} или \code{return\_void()} в зависимост от типа на резултата; и \code{unhandled\_exception()}, който запазва \code{std::exception\_ptr} в обещанието (promise) за по-късно повторно изхвърляне. При изчакване чрез \code{operator co\_await()} текущата корутина се регистрира като продължение на изчакваната задача и се възобновява автоматично, когато тази задача достигне \code{final\_suspend()}.

Помощникът \code{sync\_wait()} осигурява мост към некорутинен код. Той създава нова кратка корутина, регистрирана като продължение, и изпълнява блокираща барутера, докато тази корутина не завърши. Това е важно за единичните тестове и за гранични сценарии, в които приложението още не работи в корутинен стил.

\subsubsection{\code{ThreadPool} и универсален път на изнасяне}

\code{ThreadPool} в \path{lib/include/ORM/async/thread_pool.hpp} е стандартно реализирано обединение от ограничен брой работни нишки. Той поддържа \code{std::queue} за задачи, защитен с мутекс, и \code{std::condition\_variable} за уведомяване. Конструкторът приема брой работни нишки и стартира всяка от тях с цикъл, който изважда една задача и я изпълнява до изчерпване на опашката или сигнал за спиране.

Истинският архитектурен мост към корутините се намира в \code{run\_on\_pool()}. Тази функция връща обект, удовлетворяващ договора за изчакване (awaitable): \code{await\_ready()} винаги връща \code{false}; \code{await\_suspend(handle)} публикува извикваемия обект (callable) към обединението и съхранява \code{handle} плюс мястото за резултата; \code{await\_resume()} извлича резултата (или повдига изключение, ако работата е завършила с грешка). По този начин синхронният код на конектора, който блокира за драйверно действие, се изпълнява в работна нишка, а извикващата корутина се спира и възобновява автоматично, когато работата приключи. Този универсален път е особено важен за драйвери без неблокиращ интерфейс --- честният архитектурен избор е блокиращите извиквания да бъдат изолирани в обединението, а потребителят да получи единен интерфейс, основан на \code{Task<T>}.

\subsubsection{\code{IoContext} и интеграция с неблокиращите интерфейси на драйверите}

За драйвери с неблокиращ интерфейс библиотеката добавя реакторен слой \code{IoContext} (\path{lib/include/ORM/async/io_context.hpp}). Той предоставя \code{run()}, \code{run\_one()}, \code{stop()}, \code{post()}, \code{schedule()}, \code{watch\_readable(fd)} и \code{watch\_writable(fd)}. Системните примитиви за наблюдение зависят от платформата: на Linux се използва \code{epoll}~\cite{epoll_man} или \code{io\_uring}~\cite{io_uring}, на macOS и BSD --- \code{kqueue}~\cite{kqueue_man}, а на Windows --- I/O completion ports (IOCP)~\cite{iocp_msdn}. Извикванията \code{watch\_readable(fd)} и \code{watch\_writable(fd)} са обекти за изчакване, които не извършват входно-изходно действие; те само спират корутината и я възобновяват, когато съответният файлов описател стане готов.

Реализацията на \code{AsyncPostgreSQLDB} илюстрира как този слой се комбинира с драйвер, предоставящ неблокиращ интерфейс. Асинхронното свързване използва \code{PQconnectPoll()}, като корутината последователно изчаква \code{watch\_writable()} или \code{watch\_readable()} според състоянието на крайния автомат на \code{libpq}. При изпълнение на заявка \code{PQsendQueryParams()} стартира изпращането, \code{PQflush()} се извиква в цикъл с изчакване за готовност за запис, а \code{PQisBusy()} и \code{PQconsumeInput()} се обединяват с изчакване за готовност за четене.

\subsubsection{Обединение от връзки и транзакции при асинхронен достъп}

Помощниците \code{async\_connection\_pool<DB, N>}, \code{async\_connection\_guard<DB>} и \code{async\_transaction\_guard<DB>} разширяват синхронния обединителен слой към реализация, интегрирана с корутинния модел. Методът \code{acquire()} на обединението връща \code{Task<async\_connection\_guard<DB>>}. Вътрешно той използва \code{run\_on\_pool()} и условна променлива, за да изчака свободен слот, без извикващата корутина да бъде изложена на блокиращ интерфейс. Получената охрана следва идиомата за придобиване на ресурс при инициализация (Resource Acquisition Is Initialization, RAII): при разрушаване тя връща връзката към празно (idle) състояние и уведомява чакащите.

\code{async\_begin\_transaction()} първо изнася \code{BEGIN} към обединението и след това връща \code{async\_transaction\_guard<DB>}. Методите \code{commit()} и \code{rollback()} са \code{Task<void>} операции; деструкторът на охраната изпълнява защитно \code{ROLLBACK}, ако транзакцията не е финализирана изрично. Този дизайн е едновременно удобен (успешният път е интегриран с корутинния модел) и консервативен (пътят за почистване остава детерминистичен).

\subsection{Реализация на конекторите по системи за съхранение}
\label{subsec:backends}

Всеки конкретен конектор е специализация на \code{connector\_trait<DB>}, разположена в собствена поддиректория на \path{lib/include/ORM/db/connectors/}. Реализациите следват една и съща структура: декларация на маркерите за способности; дефиниция на \code{wire\_type} и \code{cursor\_type}; претоварвания на \code{execute(...)} за \code{select\_query}, \code{insert\_query}, \code{update\_query} и \code{delete\_query}; и при необходимост --- куки за транзакции и за миграция на схемата.

\subsubsection{SQLite}

\code{SQLiteDB} е базовият релационен конектор. Той се реализира директно върху C-интерфейса на SQLite~\cite{sqlite_api}, без външни обвивки. Връзката се представя чрез \code{sqlite3*}, а изпълнението --- чрез цикъл от \code{sqlite3\_prepare\_v2()}, \code{sqlite3\_bind\_*()}, \code{sqlite3\_step()} и \code{sqlite3\_column\_*()}. Заместителите се материализират като позиционни въпросителни знаци (\code{?}), което съответства напълно на естествения механизъм на драйвера. Конекторът декларира пълен набор от способности за релационен модел: \code{supports\_joins}, \code{supports\_transactions} и \code{supports\_aggregation}.

Изобразяването на междинното представяне към SQL се извършва чрез последователно обхождане на дървото от правила (\code{Rule}), полета (\code{field}) и заместители (\code{Placeholder}). За \code{select\_query} се изграждат фрагменти \code{SELECT}, \code{FROM}, \code{JOIN}, \code{WHERE}, \code{GROUP BY}, \code{ORDER BY} и \code{LIMIT} в строго фиксиран ред. Резултатът от изпълнението се материализира (англ. \emph{hydrate} --- утвърден термин в литературата по обектно-релационно картографиране; функцията в кода съответно носи името \code{hydrate\_result<Response>}) в \code{orm::result<Row, FieldTuple>}, където \code{Row} е tuple от извлечените колони, а \code{FieldTuple} запазва метаинформацията за проекцията. Двата термина описват едно и също действие в различен контекст: \emph{материализация} е по-широкото название от литературата по бази данни, а \emph{hydration} --- утвърденото в обектно-релационното картографиране; в текста по-нататък те се използват като взаимозаменяеми. Така достъпът до полета е възможен както чрез позиционен индекс (\code{get<I>}), така и чрез указател към член (\code{get\_field<\&T::m>}).

SQLite е и референтният сценарий за слоя на миграциите. \code{SQLiteDB::ddl\_for()} обработва \code{create\_table\_op}, \code{add\_column\_op}, \code{drop\_column\_op} и \code{alter\_column\_type\_op}, превръщайки ги в съответните DDL изрази на езика за описание на данни. Обединението с \code{migrate<DB>::diff()} дава пълноценен поток за развитие на схемата, проверен в \path{tests/unit/test_schema_migration.cpp}.

\subsubsection{PostgreSQL}

\code{PostgreSQLDB} се реализира върху \code{libpq}~\cite{libpq}. Двойната форма \code{PostgreSQLDB} и \code{PostgreSQLLiveDB} разделя двата сценария: първият е достатъчен за единични и контрактни тестове и не изисква жива (live) база данни, докато вторият работи срещу истински сървър. Заместителите се материализират като \code{\$1}, \code{\$2} и т.н., което съответства на естествения позиционен механизъм на драйвера. Конекторът декларира същия пълен набор от релационни способности като SQLite, но добавя и поддръжка за по-сложни типове и за конкурентно изпълнение чрез \code{supports\_concurrent\_execute}.

Особеност е, че PostgreSQL предоставя неблокиращ интерфейс. Това позволява \code{AsyncPostgreSQLDB} да декларира \code{supports\_async} и да реализира \code{async\_execute()} чрез \code{IoContext} и крайния автомат на \code{libpq}. На практика когато потребителят използва \code{async\_db<PostgreSQLDB>}, статичното разклонение в \code{async\_db<DB>::operator<<} избира пътя през неблокиращия интерфейс на драйвера вместо изнасяне в обединението от нишки. Интеграцията през неблокиращия интерфейс на драйвера е реализирана за повече от една система за съхранение. PostgreSQL е каноничният представителен случай, описан тук; в Раздели~\ref{subsec:backends} са разгледани и реализациите за MySQL/MariaDB (\code{AsyncMySQLDB} чрез двойките \code{\_start}/\code{\_cont}), Cassandra (\code{AsyncCassandraDB} чрез мост от \code{cass\_future\_set\_callback} към \code{std::coroutine\_handle::resume}) и Redis (\code{AsyncRedisDB} върху hiredis async API). Системите SQLite, MongoDB и Neo4j не предоставят native неблокиращ интерфейс на драйвера, поради което за тях \code{async\_db<DB>} остава върху универсалния обединителен (thread-pool) път --- честен инженерен избор, а не пропуск.

Пълноценното изпълнение на неблокиращите пътища зависи и от платформата: реакторният слой \code{IoContext} има четири конкретни реализации --- \code{kqueue\_context.cpp} (macOS, \code{EV\_ONESHOT}), \code{uring\_context.cpp} (Linux, io\_uring~\cite{io_uring} с еднократен \code{IORING\_OP\_POLL\_ADD}), \code{epoll\_context.cpp} (Linux, \code{EPOLLONESHOT} върху \code{epoll(7)}~\cite{epoll_man} --- използва се като резервен вариант, когато liburing не е наличен) и \code{iocp\_context.cpp} (Windows, \code{WSAPoll}~\cite{iocp_msdn}). Изборът между двата варианта за Linux се конфигурира чрез CMake опцията \code{ORM\_LINUX\_REACTOR} (стойности \code{auto}, \code{uring} или \code{epoll}); всички реализации споделят един и същ договор (\code{watch\_readable}, \code{watch\_writable}, \code{post}, \code{schedule}, \code{run}). При шифрована връзка през Schannel (вграденият криптографски слой на Windows) пълната асинхронна функционалност на MySQL зависи от външна поправка в библиотеката \code{MariaDB Connector/C} (\code{mariadb-connector-c\#310}~\cite{mariadb_connector_c_pr_310}, все още неприета), която отстранява два дефекта в асинхронния разпределител за шифрован вход-изход.

\begin{lstlisting}[language=C++,caption={Опростена реализация на \code{execute} за PostgreSQL},label={lst:postgresql_execute}]
template <>
struct connector_trait<PostgreSQLDB>
{
    using supports_joins              = void;
    using supports_transactions       = void;
    using supports_aggregation        = void;
    using supports_concurrent_execute = void;

    template <typename T> struct wire_type { using type = T; };
    using cursor_type = postgresql_cursor;

    template <typename Response, typename... Rest>
    static auto execute(PostgreSQLDB& db,
        select_query<Response, Rest...> q, auto&&... params)
    {
        auto [sql, n] = render_postgres_sql(q);
        const char* values[sizeof...(params)];
        bind_params(values, std::forward<decltype(params)>(params)...);

        PGresult* res = PQexecParams(db.conn(), sql.c_str(),
            n, nullptr, values, nullptr, nullptr, 0);
        return hydrate_result<Response>(res);
    }
};
\end{lstlisting}

Listing~\ref{lst:postgresql_execute} показва, че реализацията не съдържа никаква специфична за приложението логика --- тя само изобразява междинното представяне към SQL, свързва параметрите към масива от низове, изисквания от \code{PQexecParams}, и материализира резултата чрез \code{hydrate\_result<Response>}, която прилага същата проекция, запазена в типа \code{Response}.

\subsubsection{MySQL}

\code{MySQLDB} следва същата структура като PostgreSQL, но върху \code{MySQL Connector/C}~\cite{mysql_connector}. Основната разлика е в модела на свързване на параметри: MySQL използва позиционни въпросителни знаци (\code{?}), а не номерирани заместители като PostgreSQL. Това не е трудност за междинното представяне, защото операцията по изобразяване (\code{render\_mysql\_sql}) приема един и същ обект, но генерира различен текст на заявката. Реализацията използва крайния автомат на \code{MYSQL\_STMT} за подготовка, свързване и изпълнение, а резултатите се извличат чрез \code{mysql\_stmt\_fetch()}.

Маркерите за способности са идентични с тези на PostgreSQL: пълен релационен профил с поддръжка на конкурентно изпълнение. MySQL предоставя неблокиращ интерфейс под формата на двойките \code{\_start}/\code{\_cont}~\cite{mysql_capi_async} (\code{mysql\_real\_connect\_start}/\code{\_cont}, \code{mysql\_real\_query\_start}/\code{\_cont}, \code{mysql\_store\_result\_start}/\code{\_cont} при активиран \code{MYSQL\_OPT\_NONBLOCK}), и \code{AsyncMySQLDB} ги използва пряко през реакторния слой \code{IoContext}: всяко прекъсване на крайния автомат на драйвера се изразява като \code{co\_await ctx.watch\_readable(fd)} или \code{co\_await ctx.watch\_writable(fd)} според върнатата маска \code{MYSQL\_WAIT\_READ}/\code{MYSQL\_WAIT\_WRITE}. Емпиричната проверка срещу жив MariaDB~10.11 (50 последователни заявки, голям резултатен набор от 1500 реда, 10 конкурентни корутинни връзки) не показа deadlock, грешни резултати или мрежови прекъсвания. Релационната преносимост между PostgreSQL и MySQL се валидира от факта, че един и същ обект на заявката се изпълнява успешно от двата конектора, въпреки диалектните различия.

\subsubsection{MongoDB}

\code{MongoDB} е представителят на документния модел и се реализира върху \code{libmongoc}~\cite{mongoc_client_pool}. Реализацията изобразява междинното представяне към филтри и проекции в двоично представяне на JSON (Binary JSON, BSON). Полето \code{field<\&T::m>} се материализира като име на ключ в документа, а правилата за равенство, неравенство и логическо съчетаване --- като операторите \code{\$eq}, \code{\$ne}, \code{\$and} и \code{\$or}. Заявката \code{select(...).where(...)} се превръща в извикване на \code{mongoc\_collection\_find\_with\_opts()}, чийто резултат се материализира в \code{Row} чрез последователно извличане на стойностите от документа според проекцията.

\code{connector\_trait<MongoDB>} съзнателно не декларира \code{supports\_joins}, \code{supports\_transactions} или \code{supports\_aggregation} в стиловете на релационния модел. Това не е пропуск, а коректно описание на договора: документната семантика не предлага пряк еквивалент на свързване чрез външен ключ или на SQL агрегация. Отсъствието на тези маркери се проверява изрично в \path{tests/unit/test_mongodb_connector.cpp}, където опитът за заявка със свързване води до неуспех при компилация чрез \code{static\_assert}.

Релацията със стратегията \code{store\_as::embed} от обектния слой се проявява именно тук: при MongoDB вграждащият вариант се материализира като вложен документ, а препращащият (\code{store\_as::reference}) --- като отделен идентификатор, който трябва да бъде извлечен чрез следваща заявка. Това е още един пример за принципа, че логическата форма е обща, но физическата реализация зависи от системата за съхранение.

\subsubsection{Redis}

\code{RedisDB} стъпва върху ключ-стойност модела и се реализира върху \code{hiredis}~\cite{hiredis}. Обектът на ниво съхранение се картографира по два основни начина според броя на полетата. Когато обектът има едно скаларно поле (или едно сериализирано представяне), той се записва чрез \code{SET key value} и се извлича чрез \code{GET key}. Когато има няколко полета, се използват хеш-командите \code{HSET key field value} и \code{HGETALL key}, които позволяват многополева структура в рамките на едно Redis-ключово пространство. Името на логическия контейнер от \code{table\_name\_trait<T>} участва в генерирането на префикс за ключа, например \code{users:42} за обект от тип \code{User} с идентификатор 42.

Маркерите за способности са най-ограничителни в цялата библиотека: \code{connector\_trait<RedisDB>} не декларира нито един от четирите релационни маркера (свързвания, транзакции, агрегация, групиране). Това отразява честно естеството на Redis като система за достъп по ключ. Извличането може да се извършва само по първичен ключ; за по-сложни заявки потребителят трябва да поддържа допълнителни ключове за индексация, която не е автоматизирана в библиотеката. Този ограничен договор се валидира от \path{tests/unit/test_redis_connector.cpp}, където отсъствието на способностите се проверява изрично.

\subsubsection{Графов и ширококолонен сценарий (Neo4j и Cassandra)}

За пълнота на анализа е важно да се отбележи, че библиотеката съдържа и реализации на \code{Neo4jDB} (върху \code{libneo4j-client}~\cite{libneo4j_client}) и \code{CassandraDB} (върху драйвера за Cassandra с механизъм на бъдещи стойности~\cite{cassandra_futures}). Първият превръща обекти в етикетирани върхове и свойства, а заявките --- в \code{MATCH}, \code{RETURN} и \code{WHERE} фрагменти на Cypher. Вторият работи върху Cassandra Query Language (CQL) с особени ограничения, наложени от разделящия (partition) ключ. Подробно описание на тези два конектора не се включва в настоящата глава поради ограничения в обема; те се валидират в \path{tests/unit/test_neo4j_connector.cpp} и \path{tests/unit/test_cassandra_connector.cpp} съответно. Същевременно техните маркери за способности са включени в обобщителната Таблица~\ref{tab:backend_capability_matrix}.

\subsubsection{Обобщителна матрица на способностите}

Реализациите по системи за съхранение могат да бъдат сведени до компактна матрица на способностите. Тя описва архитектурната позиция на библиотеката спрямо съответната система за съхранение в текущата ревизия на хранилището, а не универсална истина за самите бази данни.

\begin{table}[H]
\centering
\small
\caption{Способности на основните системи за съхранение според специализациите и тестовете}
\label{tab:backend_capability_matrix}
\begin{tabular}{@{}L{2.0cm}L{1.1cm}L{1.5cm}L{1.3cm}L{2.0cm}L{1.9cm}L{3.5cm}@{}}
\toprule
\textbf{Backend} & \textbf{Свър.} & \textbf{Транз.} & \textbf{Агр.} & \textbf{Upsert} & \textbf{Масово} & \textbf{Коментар} \\
\midrule
SQLite & да & да & да & ограничен & ограничен & пълен релационен референтен път \\
\hdashline
PostgreSQL & да & да & да & специфичен & възможен & силен релационен договор с жива проверка \\
\hdashline
MySQL & да & да & да & специфичен & възможен & оперативен път с диалектни различия \\
\hdashline
MongoDB & не & не & не & аналог & не е централно & документен модел без релационни маркери \\
\hdashline
Redis & не & не & не & аналог & ограничен & ключ-стойност с минималистичен договор \\
\hdashline
Neo4j & аналог & да & не & не е централно & не е централно & графова система с акцент върху транзакциите \\
\hdashline
Cassandra & не & да & не & не е централно & възможен & ширококолонен с ограничен набор от способности \\
\bottomrule
\end{tabular}
\end{table}

Таблица~\ref{tab:backend_capability_matrix} показва, че контролираната нееднородност е централна за дизайна. Релационните системи за съхранение концентрират по-богат набор от способности, а документните, графовите и ключ-стойност вариантите се оценяват през други критерии. Това е инженерно честното положение --- библиотеката не обещава „универсална база данни през един интерфейс“, а \emph{унифициран логически слой с явни граници спрямо системата за съхранение}.

\subsection{Верификация и тестване}
\label{subsec:verification}

Верификационната стратегия следва архитектурата на библиотеката. Слоят, проверяван по време на компилация, се валидира чрез концепции, ограничения чрез маркери за способности и типизирано изграждане на междинното представяне. Слоят, проверяван по време на изпълнение, се покрива от единични тестове за изобразяване (rendering), за обвързване на параметри, за нишкова безопасност, за миграции, за корутинната инфраструктура и за асинхронния достъп. Най-външният слой се валидира от интеграционни тестове срещу реални системи за съхранение, така че се избягва смесване на коректността на имитиращи (mock) обекти с реалното поведение на драйверите. В текущата ревизия пълният автоматизиран набор включва 271 теста при изпълнение чрез \code{ctest}.

\subsubsection{Стратегия на верификацията}

Стратегията се основава на принципа, че всяко архитектурно твърдение трябва да има поне един пряк тестов корелат. Това включва както позитивни проверки (декларираните способности водят до успешно изпълнение на съответните операции), така и негативни проверки (отсъствието на маркер за свързване води до отказ при компилация). Негативната верификация е особено важна за конекторите MongoDB и Redis --- именно тя гарантира, че договорът на конектора отразява честно неговите ограничения, а не само неговите възможности.

\begin{table}[H]
\centering
\small
\caption{Основни източници на верификационни доказателства}
\label{tab:evidence_matrix}
\begin{tabularx}{\textwidth}{L{3.1cm}L{4.8cm}Z}
\toprule
\textbf{Източник} & \textbf{Какво валидира} & \textbf{Представителни артефакти} \\
\midrule
Ограничения при компилация и междинно представяне & типова коректност, маркери за способности и договори на конекторите & \path{tests/unit/test_query.cpp}, \path{tests/unit/test_postgresql_connector.cpp}, \path{tests/unit/test_mongodb_connector.cpp} \\
\hdashline
Асинхронна инфраструктура & \code{Task<T>}, \code{ThreadPool}, \code{IoContext}, \code{async\_db<DB>}, обединение и транзакции & \path{tests/unit/test_task.cpp}, \path{tests/unit/test_thread_pool.cpp}, \path{tests/unit/test_io_context.cpp}, \path{tests/unit/test_async_connector.cpp} \\
\hdashline
Миграции и нишкова безопасност & разлика в схемите, отмяна, охрани на връзките и конкурентна безопасност & \path{tests/unit/test_schema_migration.cpp}, \path{tests/unit/test_thread_safety.cpp} \\
\hdashline
Живи интеграции на системите за съхранение & поведение върху реални SQLite, PostgreSQL, MySQL, MongoDB и Redis среди & \path{tests/integration/test_sqlite.cpp}, \path{tests/integration/test_postgresql_live.cpp}, \path{tests/integration/test_mongodb_live.cpp}, \path{tests/integration/test_redis_live.cpp} \\
\hdashline
Емпирична оценка & пропускливост (throughput) при симулирано изчакване и изолирани измервания на режийното време (overhead) & \path{benchmarks/bench_async.cpp} \\
\bottomrule
\end{tabularx}
\end{table}

\subsubsection{Каталог на единичните и интеграционните тестове}

Единичните тестове покриват всички подсистеми на библиотеката: типовете и обектите (\path{test_property.cpp}, \path{test_relationship.cpp}); средното представяне на заявката (\path{test_query.cpp}); миграциите (\path{test_schema_migration.cpp}); нишковата безопасност (\path{test_thread_safety.cpp}); и асинхронната инфраструктура (\path{test_task.cpp}, \path{test_thread_pool.cpp}, \path{test_io_context.cpp}, \path{test_async_connector.cpp}). Допълнително \path{test_cpp26_reflection.cpp} експериментира с ранен път за интеграция с бъдещата функционалност за отражение (reflection) на C++26.

Контрактните тестове по системи за съхранение (Таблица~\ref{tab:backend_contract_tests}) валидират както позитивни, така и негативни способности. Те са разположени в \path{tests/unit/} и не изискват жива база данни --- използват имитиращи обекти и проверки чрез \code{static\_assert}, които установяват, че опитът да се използва неподдържана операция не преминава компилацията.

\begin{table}[H]
\centering
\small
\caption{Контрактни тестове по системи за съхранение}
\label{tab:backend_contract_tests}
\begin{tabularx}{\textwidth}{L{3.4cm}L{3.1cm}Z}
\toprule
\textbf{Файл} & \textbf{Фокус върху системата за съхранение} & \textbf{Тип доказателство} \\
\midrule
\path{test_postgresql_connector.cpp} & \code{PostgreSQLDB} & съвместимост с \code{is\_connector}, проверки на способностите и изобразяване на параметри \\
\hdashline
\path{test_mysql_connector.cpp} & \code{MySQLDB} & договор за свързвания и транзакции, диалектни особености и очаквания при обвързване \\
\hdashline
\path{test_mongodb_connector.cpp} & \code{MongoDB} & негативна валидация на способностите и изобразяване на филтри в BSON \\
\hdashline
\path{test_redis_connector.cpp} & \code{RedisDB} & отсъствие на свързвания, транзакции и агрегация; материализация на команди \\
\hdashline
\path{test_cassandra_connector.cpp} & \code{CassandraDB} & ограничения на способностите и съответствия с CQL \\
\hdashline
\path{test_neo4j_connector.cpp} & \code{Neo4jDB} & транзакционна способност и графова материализация на заявки \\
\bottomrule
\end{tabularx}
\end{table}

\subsubsection{Интеграционни и реално свързани с живи системи за съхранение тестове}

Интеграционните тестове са разделени на два слоя. Първият е \path{tests/integration/test_mockdb.cpp} --- винаги наличен интеграционен базов сценарий, който проверява съставянето на заявки, обвързването на параметрите, подготвените заявки и CRUD сценариите без външна система за съхранение. Той е особено полезен, защото гарантира, че логическият слой е стабилен независимо от наличието на жива база данни. Вторият слой включва тестовете срещу реални системи за съхранение, които се активират условно: \path{test_sqlite.cpp} (локален интеграционен тест, винаги наличен), \path{test_postgresql_live.cpp}, \path{test_mysql_live.cpp}, \path{test_mongodb_live.cpp}, \path{test_redis_live.cpp}, \path{test_neo4j_live.cpp} и \path{test_cassandra_live.cpp}.

Условното активиране се извършва чрез променливи на обкръжението: ако съответната променлива (например \code{ORM\_TEST\_POSTGRESQL=1}) е зададена, тестът се изпълнява срещу инстанция, предоставена от Docker-базирана инфраструктура. Това позволява разработчикът да изпълнява пълния набор от тестове в среда без външни зависимости, а екипите за качество да активират по-задълбочените проверки в специално подготвена среда. Разделението избягва смесване на коректността на имитиращи обекти с действителното поведение на драйверите.

\subsubsection{Емпирични резултати}

Емпиричната оценка използва библиотеката Google Benchmark~\cite{google_benchmark_user_guide} и е реализирана в \path{benchmarks/bench_async.cpp}. Тя не симулира пълния жизнен цикъл на конкретна база данни, а изолира архитектурния въпрос дали корутинно базираното изпълнение освобождава работните нишки по-ефективно от синхронното блокиране, когато между две кратки изчислителни фази съществува интервал на изчакване. Наборът е разделен на две групи. Първата измерва пропускливост (throughput) на цели партиди от задачи: синхронният базов сценарий блокира работна нишка, докато асинхронният използва \code{SimulatedAsyncIO} обект за изчакване, който спира корутината, планира отложено възобновяване чрез таймерна опашка и връща продължението обратно към обединението от нишки. Втората група измерва изолирани режийни времена: създаване и завършване на \code{Task<int>}, обхождане в две посоки на \code{ThreadPool::post()} и вложени вериги от \code{co\_await}.

\begin{table}[H]
\centering
\small
\caption{Представителни резултати в режим release}
\label{tab:throughput_results}
\begin{tabular}{L{1.3cm}L{1.5cm}L{1.9cm}L{2.4cm}L{2.4cm}L{1.7cm}}
\toprule
\textbf{Нишки} & \textbf{Задачи} & \textbf{Изчакване} & \makecell[tl]{\textbf{Синхронна}\\\textbf{пропускливост}} & \makecell[tl]{\textbf{Асинхронна}\\\textbf{пропускливост}} & \textbf{Ускорение} \\
\midrule
4 & 100 & 10~$\mu$s & 194{,}6\,k оп/с & 658{,}5\,k оп/с & 3.38$\times$ \\
\hdashline
4 & 100 & 100~$\mu$s & 29{,}8\,k оп/с & 458{,}2\,k оп/с & 15.35$\times$ \\
\hdashline
4 & 100 & 1000~$\mu$s & 3{,}22\,k оп/с & 75{,}28\,k оп/с & 23.36$\times$ \\
\hdashline
4 & 1000 & 100~$\mu$s & 30{,}08\,k оп/с & 670{,}5\,k оп/с & 22.29$\times$ \\
\hdashline
8 & 100 & 100~$\mu$s & 57{,}13\,k оп/с & 258{,}5\,k оп/с & 4.53$\times$ \\
\bottomrule
\end{tabular}
\end{table}

Най-важният извод от Таблица~\ref{tab:throughput_results} е, че предимството на корутинния модел нараства не когато локалната изчислителна работа се увеличи, а когато фазата на изчакване започне да доминира над нея. При 4 работни нишки и 100 задачи разликата между 10 и 1000 микросекунди изчакване увеличава относителното предимство на асинхронния вариант от 3.38$\times$ до 23.36$\times$. Това е точно очакваното поведение: при синхронния базов сценарий нишките се изчерпват от блокиране, докато при корутинния модел те се освобождават за друга работа. Вторият извод е, че ефектът се запазва и при по-големи партиди --- при 4 нишки, 1000 задачи и 100~$\mu$s изчакване асинхронният вариант достига 670{,}5\,k операции в секунда срещу 30{,}08\,k при синхронния.

\begin{figure}[H]
\centering
\begin{tikzpicture}[x=1.95cm,y=0.24cm]
    \draw[->, thick] (0,0) -- (0,26) node[above] {\small ускорение};
    \draw[->, thick] (0,0) -- (6.4,0);

    \foreach \y in {5,10,15,20,25}
    {
        \draw[gray!35] (0,\y) -- (6.2,\y);
        \node[left] at (0,\y) {\scriptsize \y};
    }

    \foreach \x/\h/\labeltext/\valuetext in {
        0.7/3.38/{4 нишки\\100 задачи\\10~$\mu$s}/3.38,
        1.9/15.35/{4 нишки\\100 задачи\\100~$\mu$s}/15.35,
        3.1/23.36/{4 нишки\\100 задачи\\1000~$\mu$s}/23.36,
        4.3/22.29/{4 нишки\\1000 задачи\\100~$\mu$s}/22.29,
        5.5/4.53/{8 нишки\\100 задачи\\100~$\mu$s}/4.53}
    {
        \fill[green!28, draw=black] (\x-0.27,0) rectangle (\x+0.27,\h);
        \node[font=\scriptsize, align=center, above] at (\x,\h) {\valuetext$\times$};
        \node[font=\scriptsize, align=center, below=7pt] at (\x,0) {\labeltext};
    }
\end{tikzpicture}
\caption{Относително ускорение на асинхронната пропускливост спрямо синхронната}
\label{fig:async_speedup}
\end{figure}

При нулева изкуствена фаза на изчакване (Таблица~\ref{tab:zero_latency_results}) ситуацията се променя. Пропускливостта на асинхронния вариант е сравнима или дори малко по-ниска от синхронната. Това е важен коректив: корутинният слой не „създава“ пропускливост сам по себе си --- той добавя малко, но измеримо допълнително време, което може да бъде компенсирано само частично, когато няма реално изчакване.

\begin{table}[H]
\centering
\small
\caption{Представителни резултати при нулева фаза на изчакване}
\label{tab:zero_latency_results}
\begin{tabular}{L{1.4cm}L{1.5cm}L{2.8cm}L{2.9cm}L{3.0cm}}
\toprule
\textbf{Нишки} & \textbf{Задачи} & \makecell[tl]{\textbf{Синхронна}\\\textbf{пропускливост}} & \makecell[tl]{\textbf{Асинхронна}\\\textbf{пропускливост}} & \makecell[tl]{\textbf{Отношение}\\\textbf{асинх./синхр.}} \\
\midrule
4 & 100 & 841{,}0\,k оп/с & 989{,}4\,k оп/с & 1.18$\times$ \\
\hdashline
4 & 1000 & 1{,}232\,M оп/с & 1{,}184\,M оп/с & 0.96$\times$ \\
\hdashline
8 & 100 & 810{,}9\,k оп/с & 514{,}4\,k оп/с & 0.63$\times$ \\
\bottomrule
\end{tabular}
\end{table}

Изолираните измервания показват, че режийното време на корутинната инфраструктура не доминира в сценариите, ограничени от вход-изход. В режим release \code{BM\_TaskCreateAndWait} измерва приблизително 71{,}6\,ns за създаване и завършване на тривиална задача, \code{BM\_ThreadPoolPost} остава около 5{,}0~$\mu$s независимо от размера на обединението от нишки, а \code{BM\_NestedCoAwait} достига около 324\,ns при дълбочина 10 и около 2{,}375~$\mu$s при дълбочина 100. Тези числа са с порядъци по-малки от 100--1000~$\mu$s фази на изчакване, при които асинхронната пропускливост печели най-много. С други думи, наблюдаваното ускорение не е артефакт на измервателния инструмент, а следствие от това, че работните нишки престават да бъдат заключени в блокиращо изчакване.

\subsubsection{Покритие и ограничения на верификацията}

Емпиричната оценка съзнателно използва симулирана фаза на изчакване, а не пълно изпълнение от край до край срещу мрежов сървър за база данни. Това е ограничение, ако целта е да се прогнозира абсолютна пропускливост за реална инсталация на PostgreSQL или Cassandra, но е методологично предимство, ако целта е да се изолира моделът на изпълнение от влиянието на планировчика на заявките, мрежовите колебания, кеширането в драйвера и спецификите на сторидж двигателя. Следователно представените резултати не доказват, че „асинхронното е винаги по-бързо“; те доказват по-тясната, но по-научно защитима теза, че когато архитектурата работи в режим, ограничен от вход-изход, и фазата на изчакване доминира над локалната изчислителна работа, корутинно базираният модел позволява съществено по-добро използване на работните ресурси от синхронния базов сценарий.

Тестовата верификация и емпиричната оценка покриват двата основни въпроса на дипломната работа. От една страна библиотеката показва формална и тестово подкрепена коректност. От друга страна асинхронният изпълнителен модел показва измерима практическа полза точно в сценария, за който е проектиран. Така емпиричната оценка не стои встрани от архитектурата, а потвърждава нейния централен инженерен и научен смисъл.

\begin{table}[H]
\centering
\small
\caption{Връзка между целите на разработката и постигнатите резултати}
\label{tab:goal_result_mapping}
\begin{tabularx}{\textwidth}{L{3.0cm}L{3.5cm}Y}
\toprule
\textbf{Цел} & \textbf{Реализиран механизъм} & \textbf{Доказателствена опора} \\
\midrule
Типизиран модел на данните & \code{property}, \code{relationship}, \code{table\_name\_trait} & заглавните файлове на обектния слой; единичните тестове за свойства и връзки \\
\hdashline
Типизиран слой на заявките & съставители и междинно представяне & файловете на слоя на заявките, \path{tests/unit/test_query.cpp}, \path{tests/integration/test_mockdb.cpp} \\
\hdashline
Разширяем слой за много системи за съхранение & договор на конектора и маркери за способности & заглавните файлове на конектора и тестовете по системи за съхранение \\
\hdashline
Подход, работещ със схемата & \code{migrate<DB>} и куки за DDL & \path{tests/unit/test_schema_migration.cpp} \\
\hdashline
Безопасна оперативна употреба & помощници за нишкова безопасност и транзакции & \path{tests/unit/test_thread_safety.cpp} \\
\hdashline
Асинхронно изпълнение & \code{Task<T>}, \code{ThreadPool}, \code{IoContext}, \code{async\_db<DB>} & асинхронните единични тестове и \path{benchmarks/bench_async.cpp} \\
\bottomrule
\end{tabularx}
\end{table}

Таблица~\ref{tab:goal_result_mapping} показва, че работата не остава на равнището на декларативен проектен документ: за всяка основна цел има конкретен реализиран механизъм и поне един пряк технически корелат в хранилището. Практическата стойност на това съответствие е, че библиотеката предоставя стабилна основа за по-нататъшно инженерно развитие --- обектното описание, междинното представяне на заявката, договорът на конектора и верификационната стратегия са достатъчно ясно формулирани, за да позволят дисциплинирано добавяне на нови системи за съхранение и оперативни възможности.

\subsection{Приоритетни направления за бъдещо развитие}
\label{subsec:future_directions}

Разработените механизми очертават естествени посоки за по-нататъшно надграждане. Първо, по-голяма оперативна зрелост на вече наличните системи за съхранение --- по-широко покриване на типовете, диагностичните пътища при грешка, политиките за жизнен цикъл на връзките и поведението при по-сложни заявки. Второ, по-пълна автоматизация на схемата чрез извличане на живи метаданни от драйверите, генериране на безопасни DDL операции и анализ на съвместимостта. Трето, разширяване на интеграцията през неблокиращите интерфейси на драйверите и за други системи за съхранение (MySQL чрез \code{\_start}/\code{\_cont} двойките, Redis чрез \code{hiredis} с обратни извиквания), което би доближило библиотеката до по-нисък разход в реални мрежови сценарии. Четвърто, интегриране на бъдещата функционалност за отражение (reflection) на C++26, която би позволила още по-кратко описание на обектите без обвивката \code{property<T, Name>}. Пето, по-фина таксономия на способностите --- разграничаване между безопасно паралелно четене и пълно конкурентно изпълнение, между специфични форми на вграждане в документ или между различни нива на поддръжка на стрийминг резултати.

\begin{table}[H]
\centering
\small
\caption{Приоритетни направления за бъдещо развитие}
\label{tab:future_work_priorities}
\begin{tabularx}{\textwidth}{L{2.8cm}L{3.7cm}L{2.9cm}Y}
\toprule
\textbf{Направление} & \textbf{Засегнати подсистеми} & \textbf{Очакван ефект} & \textbf{Предизвикателство} \\
\midrule
По-зряло покритие на живи системи за съхранение & конектори, интеграционни тестове, пътища при грешка & по-висока практическа надеждност & различна оперативна семантика между системите за съхранение \\
\hdashline
Автоматизация на работата със схемата & \code{migrate<DB>}, куки за DDL, метаданни & по-близка до производствено приложение интеграция & съчетаване на компилационна и изпълнителна истина \\
\hdashline
Неблокираща интеграция чрез интерфейсите на драйверите & реакторен слой и крайни автомати на драйверите & по-нисък разход в реални мрежови сценарии & различни модели на готовност при различните драйвери \\
\hdashline
Моделиране, базирано на отражение (reflection) & обектен слой, извличане на метаданни, ергономия на интерфейса & по-малко шаблонен код и подход „модел напред“ & баланс между автоматизация и прозрачност \\
\hdashline
По-фина таксономия на способностите & \code{connector\_trait<DB>} и статични проверки & по-точен договор за нови системи за съхранение & избягване на декоративни маркери \\
\hdashline
По-богата емпирична оценка & асинхронен слой и подготвени заявки & по-точен модел на разходите & необходимост от стабилни производствени среди \\
\bottomrule
\end{tabularx}
\end{table}

Накрая, една бъдеща версия на работата следва да включва по-систематичен набор от измервания, в който се сравняват не само различните системи за съхранение, но и различни режими на изпълнение в рамките на самата библиотека --- пряко изпълнение, повторна употреба на подготвени заявки, сценарии за масово вмъкване и различни форми на материализация на резултата. Такива измервания биха показали реалната цена на различните архитектурни избори и биха обогатили както инженерната, така и научната оценка на разработката.
